% Full proofs for the separately authored represented-space foundation.

\begin{annexitem}{annex:p04-rs-finite-use}{thm:p04-rs-finite-use}{Finite use, continuity and composition}
\begin{proof}
Let \(P\) be a machine's productive domain and let \(p\in P\).
For each \(n\), its first \(n\) complete output entries are issued
after a finite number \(t\) of steps. Only finitely many input symbols
have then been inspected. Take a prefix containing all of them,
including the delimiter of each relevant Baire entry.
Every other input with that prefix follows the identical computation
through step \(t\), and therefore issues the same \(n\) entries.
Restricting to \(P\), the inverse image of each output cylinder is
open. Thus the partial output operator is continuous on \(P\).
A machine realizer is consequently a continuous realizer.

For represented maps \(f:X\rightharpoonup Y\) and
\(g:Y\rightharpoonup Z\), let \(F,G\) be their machine realizers.
On a name of \(x\) in the domain of \(g\circ f\), the first machine
is productive and produces a valid name of \(f(x)\); the second
is productive on that name and names \(g(f(x))\).
Build one machine by keeping the two simulations and a buffer.
Whenever the second simulation requests the next intermediate
symbol, advance the first until that symbol is issued, cache it,
and resume the second. Work already done and input already read
are retained. Every such request on a promised input receives its
answer after finitely many steps, and each requested output prefix
is eventually completed. The resulting stream is a name of the
composed value. Extra productive inputs outside the represented
map's domain do not violate the realizer definition.

For continuous realizers use \(G\circ F\) on
\(\{p\in\dom(F):F(p)\in\dom(G)\}\). Composition is continuous
in the subspace topology and contains every required input name.
This proves the analogous continuous assertion.
\end{proof}
\end{annexitem}

\begin{annexitem}{annex:p04-rs-invariance}{thm:p04-rs-invariance}{Both translations and an effectivity counterexample}
\begin{proof}
Write \(\delta_X,\delta_Y\) for the old representations and
\(\gamma_X,\gamma_Y\) for the equivalent replacements.
Let \(A\) translate \(\gamma_X\)-names to \(\delta_X\)-names and
\(B\) translate \(\delta_Y\)-names to \(\gamma_Y\)-names.
If \(F\) realizes \(f\) in the old representations, then
\(B\circ F\circ A\) realizes it in the new ones: on every
promised \(\gamma_X\)-name \(p\),
\[
 \gamma_Y(B(F(A(p))))
 =\delta_Y(F(A(p)))
 =f(\delta_X(A(p)))
 =f(\gamma_X(p)).
\]
Each intermediate value is defined and a valid name on that
domain. The finite-use composition construction supplies a
machine when the three operators are computable.
Use the reverse input/output translators to prove the converse.
Translating a computable point-name proves point invariance in
each direction. The same composition equations with continuous
operators prove invariance of continuous realizability.

Continuous equivalence does not suffice for the computational
claim. Finite programs form a countable set, so their productive
binary outputs form a countable subset of \(\{0,1\}^{\mathbb N}\).
The latter is uncountable: from a purported list \(b_0,b_1,\ldots\),
the stream \(a(n)=1-b_n(n)\) differs from its \(n\)-th member.
Choose a binary stream \(a\) with no computable output program.
Represent \(X=\{0,1\}^{\mathbb N}\) by identity on binary streams,
and also by
\[
 \gamma(p)(n)=p(n)\mathbin{\oplus}a(n)
 \qquad(p\in\{0,1\}^{\mathbb N}),
\]
where \(\oplus\) is exclusive-or on bits. Both translations are
the same involution \(p\mapsto p\oplus a\). They are continuous:
an \(n\)-entry output prefix depends on just the first \(n\)
input entries. Identity makes the zero stream a computable point,
whereas its unique \(\gamma\)-name is the noncomputable stream
\(a\). Therefore continuous equivalence can change computable
points, even in this simple infinite space.
\end{proof}
\end{annexitem}

\begin{annexitem}{annex:p04-rs-standard}{thm:p04-rs-standard}{Neighbourhood names and their universal property}
\begin{proof}
For each \(x\in X\), its set \(I_x=\{i:x\in U_i\}\) has an
enumeration as a set of natural numbers, without an effectivity
claim. List its members, allow repetition and fill unused places
with zero; if it is finite, use zero padding after the list.
This gives a name. If \(x\ne y\), the \(T_0\) property gives an
open set containing one and not the other. A basic set inside
that open set distinguishes \(I_x\) and \(I_y\). Thus the
representation is well defined and surjective. The empty-space
case has empty valid-name domain and satisfies the claims
vacuously.

On its domain \(E\),
\[
 \delta_{\mathrm{nb}}^{-1}(U_i)
 =\{p\in E:\exists n\ p(n)=i+1\},
\]
a union of cylinders intersected with \(E\).
Preimages of all open sets are consequently open, so the
representation is continuous.

Let \(h:D\to X\) be continuous for any \(D\subseteq\mathcal B\).
Fix an enumeration of all pairs \((i,s)\), where \(s\) is a finite
natural sequence. These pairs are countable; for example enumerate
successively by bounds on \(i\), sequence length and its entries.
For the pair at position \(j\), fix the truth value
\[
 a_j=1\quad\text{iff}\quad D\cap[s]\subseteq h^{-1}(U_i).
\]
The sequence of truth values is not assumed computable.
For \(p\in D\) define
\[
 H(p)(j)=
 \begin{cases}
 i+1,&a_j=1\text{ and }p\in[s],\\
 0,&\text{otherwise}.
 \end{cases}
\]
Each output coordinate depends on a finite prefix of \(p\);
the fixed \(a_j\) is a constant for that coordinate. Cylinders
are clopen, so \(H\) is continuous as a map to \(\mathcal B\).
Every emitted positive index names a basic neighbourhood of
\(h(p)\). Conversely, if \(h(p)\in U_i\), continuity makes
\(h^{-1}(U_i)\) an open neighbourhood of \(p\) in \(D\).
Some cylinder \([s]\) containing \(p\) has
\(D\cap[s]\subseteq h^{-1}(U_i)\).
Its pair occurs in the enumeration and emits \(i+1\).
Hence \(H(p)\) enumerates exactly \(I_{h(p)}\), lies in \(E\),
and satisfies \(\delta_{\mathrm{nb}}H=h\).
Continuity into \(E\) follows from the subspace topology.
This proves admissibility with no effective claim about the
arbitrary truth-value sequence.

To prove the quotient property, suppose
\(\delta_{\mathrm{nb}}^{-1}(V)\) is open in \(E\).
For \(x\in V\), choose a name \(p\) and a finite prefix \(s\)
such that \(E\cap[s]\subseteq\delta_{\mathrm{nb}}^{-1}(V)\).
Let \(J\) be the finite set of positive indices listed in \(s\).
The finite intersection \(W=\bigcap_{i\in J}U_i\), with empty
intersection interpreted as \(X\), is an open neighbourhood
of \(x\). Every \(y\in W\) has a name beginning with \(s\):
retain that prefix, then enumerate all of \(I_y\), with padding.
All indices in the prefix are permitted for \(y\).
Consequently \(y\in V\), so \(W\subseteq V\).
Thus \(V\) is open. Together with continuity this is the
quotient property. With target \(A\subseteq X\), use exactly
the same argument for \(x\in A\cap V\) and \(W\cap A\);
extensions name only \(y\in A\). This proves the restricted
quotient assertion for every subspace \(A\).

Finally let \(\delta:E'\to X\) be another continuous
representation. The just-proved universal property supplies a
continuous translator \(\delta\leq_{\mathrm t}\delta_{\mathrm{nb}}\).
If \(\delta\) is admissible, apply its lifting property to
\(\delta_{\mathrm{nb}}:E\to X\) to get the opposite translation.
Conversely, if the two representations are continuously
equivalent, lift any \(h:D\to X\) first to neighbourhood names
and then translate those names into \(\delta\)-names.
The composed lift is continuous and has the required equation.
This proves both directions of the characterization.
\end{proof}
\end{annexitem}

\begin{annexitem}{annex:p04-rs-topological}{thm:p04-rs-topological}{Lifting and quotient descent}
\begin{proof}
If \(f:X\to Y\) is continuous, then
\(h=f\circ\delta_X:E_X\to Y\) is continuous.
Admissibility of \(\delta_Y\), applied to the Baire subspace
\(E_X\), provides a continuous \(H:E_X\to E_Y\) with
\(\delta_YH=f\delta_X\).
Taking \(H\) as a partial operator on \(\mathcal B\) gives
a continuous realizer on all valid input names.

Conversely, restrict a continuous realizer \(G\) to \(E_X\).
For any open \(V\subseteq Y\), continuity of \(\delta_YG\)
makes
\[
 \delta_X^{-1}(f^{-1}(V))
 =(\delta_YG)^{-1}(V)
\]
open in \(E_X\). Since \(\delta_X\) is quotient,
\(f^{-1}(V)\) is open in \(X\). Hence \(f\) is continuous.
This direction uses quotient descent and continuity of the
output representation, not its admissibility.

For a partial map let \(A=\dom(f)\) and
\(E_A=\delta_X^{-1}(A)\). The restricted representation
\(\delta_A:E_A\to A\) is surjective and continuous.
Under the stated quotient hypothesis the two arguments apply
to the total map \(f:A\to Y\). The preceding annex item proves
that hypothesis for the restricted neighbourhood representation.
No general assertion that restrictions of arbitrary quotient
maps are quotient is used.
\end{proof}
\end{annexitem}

\begin{annexitem}{annex:p04-rs-cauchy-bridge}{thm:p04-rs-cauchy-bridge}{All four naming clauses and both directions}
\begin{proof}
The finite rational coding used by the translators is effective.
For a fraction \(z/d\) with \(d>0\), encode the pair
\((\zeta^{-1}(z),d-1)\), where
\(\zeta^{-1}(z)=2z\) for \(z\geq0\) and
\(\zeta^{-1}(z)=-2z-1\) otherwise.
Finite fraction arithmetic uses integer operations and
cross-multiplication with positive denominators for order tests.
If reduced syntax is required by a source fraction notation,
Euclid's remainder loop computes it: each positive remainder
is smaller than the previous divisor, so the loop terminates;
the common divisors of \((a,b)\) equal those of
\((b,a\bmod b)\), so the final nonzero divisor is their gcd.
Divide numerator and denominator by that value and retain
the sign in the numerator. This supplies the finite syntax
conversions used below without assuming a rational-name oracle.

If one decoded sequence meets the error clause for both \(x\)
and \(y\), then \(\lvert x-y\rvert\leq2^{1-n}\) for every \(n\),
so \(x=y\). For each real \(x\), rational density gives a rational
within \(2^{-n}\) at every position \(n\); choosing the first
such rational in the fixed enumeration defines a name
set-theoretically. Thus \(\delta_{\mathrm C}\) is a partial
surjection, without claiming that this selection is computable
from an unnamed \(x\).

For \(\delta_{\mathrm C}\leq_{\mathrm c}\delta_{\mathrm I}\),
enumerate all pairs \((i,k)\) effectively.
At the pair with \(I_i=(a_i,b_i)\), read \(q_k\) and emit
\(i+1\) exactly when
\[
 a_i<q_k-2^{-k}
 \quad\text{and}\quad
 q_k+2^{-k}<b_i;
\]
otherwise emit padding zero. These are decidable rational
comparisons. An emitted index contains the named \(x\).
Conversely, if \(a_i<x<b_i\), choose \(k\) with
\(2^{1-k}<\min(x-a_i,b_i-x)\). Since
\(\lvert q_k-x\rvert\leq2^{-k}\), both displayed strict
inequalities hold. Hence every interval containing \(x\)
is listed, and no other one is listed.
Each coordinate uses one finite requested input entry.

For \(\delta_{\mathrm I}\leq_{\mathrm c}\delta_{\mathrm C}\),
at output position \(n\), search the received interval indices,
retaining a finite cache, for an interval \((a,b)\) of width
less than \(2^{-n}\). Then output the rational \((a+b)/2\).
Every valid name lists arbitrarily small rational intervals
containing \(x\), so every search terminates.
Since \(a<x<b\),
\(\lvert(a+b)/2-x\rvert<(b-a)/2<2^{-n}\).
Thus the output satisfies the host clause at every position.

For a source pairwise name, completeness of the reals gives
its limit \(x\); taking \(j\to\infty\) in
\(\lvert u_j-u_n\rvert<2^{-n}\) gives
\(\lvert x-u_n\rvert\leq2^{-n}\).
The same decoded sequence is therefore a host name.
Conversely, from a host name use \(u_n=q_{n+2}\).
For \(j>n\),
\[
 \lvert u_j-u_n\rvert
 \leq2^{-j-2}+2^{-n-2}<2^{-n},
\]
and \(u_n\to x\). The shifted sequence meets the strict
source pairwise clause. Reading or writing each finite
fraction block, or its natural rational code, is a terminating
effective syntax conversion; the delimiter is not a new
mathematical datum.

For \(\delta_{\mathrm C}\leq_{\mathrm c}\rho_{\mathrm{bounds}}\),
output
\(\ell_n=q_n-2^{-n}\) and \(v_n=q_n+2^{-n}\).
All lower bounds are at most \(x\), all upper bounds at least
\(x\), and both sequences tend to \(x\).
It follows that \(\sup\ell_n=\inf v_n=x\).
Conversely, every source bound pair has
\(\ell_i\leq x\leq v_j\) for all \(i,j\).
For output \(n\), search pairs \((i,j)\) for
\(0\leq v_j-\ell_i<2^{-n}\) and output their midpoint.
The supremum and infimum premises guarantee a pair with
arbitrarily small width: choose a lower bound within half
the requested width below \(x\) and an upper one similarly
above \(x\). The search is effective over rational input
entries and terminates on every valid name.
Its midpoint is within \(2^{-n-1}\) of \(x\).
Equality of the bounds is harmless: it returns the exact
rational value. No monotonicity has been assumed.
All four representations are now connected in both
computable directions through \(\delta_{\mathrm C}\).

Open rational intervals form a countable basis of the usual
\(T_0\) real topology, so \(\delta_{\mathrm I}\) is the
neighbourhood representation covered by
\cref{thm:p04-rs-standard}.
The computable translators are continuous by finite use.
The equation \(\delta_{\mathrm C}
=\delta_{\mathrm I}\circ T_{\mathrm{CI}}\) proves continuity
of \(\delta_{\mathrm C}\).
For any continuous \(h:D\to\mathbb R\), first lift it to
interval names, then apply \(T_{\mathrm{IC}}\); the
composition is a continuous Cauchy-name lift.
This proves admissibility.
If \(\delta_{\mathrm C}^{-1}(V)\) is open in its domain,
then
\(\delta_{\mathrm I}^{-1}(V)
=T_{\mathrm{IC}}^{-1}(\delta_{\mathrm C}^{-1}(V))\)
is open. Quotientness of \(\delta_{\mathrm I}\) makes \(V\)
open in \(\mathbb R\). The converse follows from continuity.
For any subspace \(A\), both translators preserve the named
value and restrict to the names of \(A\).
The same argument using the already proved restricted
interval quotient proves the Cauchy restriction is quotient.
\end{proof}
\end{annexitem}

\begin{annexitem}{annex:p04-rs-arithmetic}{thm:p04-rs-arithmetic}{Effective precision choices and output-name verification}
\begin{proof}
Let \(q_n,r_n\) be valid names of \(x,y\).
For addition,
\[
 \lvert(q_{n+1}+r_{n+1})-(x+y)\rvert
 \leq2^{-n-1}+2^{-n-1}=2^{-n}.
\]
Negation has the same input error. The reverse triangle
inequality gives
\(\bigl|\lvert q_n\rvert-\lvert x\rvert\bigr|
\leq\lvert q_n-x\rvert\leq2^{-n}\), proving the absolute-value
construction. Subtraction follows by addition after negation.

Put \(\epsilon=2^{-n}\).
The inequalities \(q_n\leq x+\epsilon\) and
\(r_n\leq y+\epsilon\) imply
\(\max(q_n,r_n)\leq\max(x,y)+\epsilon\).
The inequalities in the other direction imply
\(\max(x,y)\leq\max(q_n,r_n)+\epsilon\).
Thus \(\max(q_n,r_n)\) is a valid approximation.
Negating the maximum of \(-q_n,-r_n\) proves the minimum
construction with the same bound.

For multiplication the zeroth entries give
\(\lvert x\rvert\leq B_x=\lvert q_0\rvert+1\) and
\(\lvert y\rvert\leq B_y=\lvert r_0\rvert+1\).
For natural \(m\), \(\lvert q_m\rvert\leq B_x+1\).
Consequently
\[
 \begin{split}
 \lvert q_mr_m-xy\rvert
 &\leq \lvert q_m\rvert\,\lvert r_m-y\rvert
       +\lvert y\rvert\,\lvert q_m-x\rvert\\
 &\leq(B_x+B_y+1)2^{-m}.
 \end{split}
\]
Search \(m=0,1,\ldots\) until the stated rational bound is
at most \(2^{-n}\). The coefficient is a known positive
rational and \(2^{-m}\to0\), so this search terminates.
Reading the requested entries and multiplying their rational
values gives the output code at position \(n\).

Each algorithm uses only finitely many input entries for
each finite output stage, performs effective finite rational
operations, and satisfies the output error clause.
Splitting an interleaved pair-stream needs a finite read to
reach each requested component entry; caches allow reuse.
Hence these are genuine productive Type-Two realizers on
all promised names. Although the precision choices can
differ between names of one input, the proved bounds force
the same decoded operation value.
\end{proof}
\end{annexitem}

\begin{annexitem}{annex:p04-rs-reciprocal}{thm:p04-rs-reciprocal}{Host error-to-limit and source strict-pairwise reciprocal proofs}
\begin{proof}
First consider the host clause. If \(x\ne0\), then
\[
 \lvert q_k\rvert-2^{-k}
 \geq\lvert x\rvert-2^{1-k}>0
\]
for all sufficiently large \(k\). Thus the certificate search
terminates on every name of a nonzero real.
For a valid zero name,
\(\lvert q_k\rvert\leq2^{-k}\) at every \(k\), so it never
passes the test and the algorithm issues no output.
When it passes, the positive rational \(L\) satisfies
\(\lvert x\rvert\geq L\).

Both precision inequalities can be checked exactly and
eventually hold as \(m\to\infty\).
For the selected \(m\),
\[
 \lvert q_m\rvert
 \geq\lvert x\rvert-2^{-m}\geq L/2>0.
\]
The rational inverse is therefore defined and
\[
 \left|\frac1{q_m}-\frac1x\right|
 =\frac{\lvert q_m-x\rvert}{\lvert q_m\rvert\lvert x\rvert}
 \leq\frac{2^{1-m}}{L^2}\leq2^{-n}.
\]
This verifies every host output position and productivity
on all nonzero names. The positive \(L\) was derived from
the name, not assumed as an additional promise.

Now retain the source pairwise clause separately.
It entails \(\lvert u_i-x\rvert\leq2^{-i}\), so when
\(x\ne0\), eventually
\(\lvert u_N\rvert>2\cdot2^{-N}\).
After the first such \(N\), for every \(i\geq N\),
\[
 \lvert u_i\rvert
 \geq\lvert u_N\rvert-\lvert u_i-u_N\rvert>2^{-N};
\]
for \(i=N\) this follows directly from the certificate.
Also \(\lvert x\rvert
\geq\lvert u_N\rvert-2^{-N}>2^{-N}\).
Thus every denominator at index \(2N+n\) is nonzero.
Let \(w_n=1/u_{2N+n}\).
For \(j>n\), the original strict pairwise premise gives
\[
 \lvert w_j-w_n\rvert
 =\frac{\lvert u_{2N+j}-u_{2N+n}\rvert}
        {\lvert u_{2N+j}\rvert\lvert u_{2N+n}\rvert}
 <2^{-2N-n}\,2^{2N}=2^{-n}.
\]
Moreover
\[
 \left|w_n-\frac1x\right|
 \leq 2^{-2N-n}\,2^{2N}=2^{-n},
\]
using the error-to-limit consequence and the same denominator
lower bounds. Therefore \(w_n\to1/x\) and the strict pairwise
output condition holds. At zero the source certificate cannot
pass either. This completes the distinct 1995-clause algorithm,
including its domain and output conditions.

For division, split the pair name into \(x,y\), run the
reciprocal algorithm on \(y\), and feed the resulting name
with the retained name of \(x\) into multiplication.
The finite-use composition construction is productive for
all \(y\ne0\) and names \(x/y\). Its promise domain is exactly
the stated mathematical domain; no membership decider was
used.
\end{proof}
\end{annexitem}

\begin{annexitem}{annex:p04-rs-zero}{thm:p04-rs-zero}{Finite-prefix obstruction and strict-separation searches}
\begin{proof}
The constant zero-code stream \(p\) names zero.
If a continuous realizer \(G\) of \(Z\) existed, its first
output entry at \(p\) would be one.
Continuity gives a finite prefix of \(k\) zero entries such
that every stream in \(\dom(G)\) with that prefix also has
first output entry one.
Put \(x=2^{-k-1}>0\).
It has a valid host name \(q\) whose first \(k\) decoded
rationals are zero and whose subsequent rationals equal
\(x\): for \(j<k\), \(x\leq2^{-j}\), and afterwards the
error is zero. All the first \(k\) rational codes can
be the zero code, just as in \(p\).
Since \(q\) is a valid name of an input in the total domain,
it lies in \(\dom(G)\). The first output entry is then one,
contradicting \(Z(x)=0\).
This proves failure of continuous realizability, and the
finite-use theorem rules out a computable realizer as well.
The argument includes \(k=0\), when the common prefix is empty.

A positive recognizer reads successive \(q_k\) and halts
as soon as \(q_k-2^{-k}>0\).
This certificate implies \(x>0\).
If \(x>0\), choose \(k\) with \(2^{1-k}<x\);
then \(q_k-2^{-k}\geq x-2^{1-k}>0\), so the recognizer
eventually halts.
The negative recognizer tests \(q_k+2^{-k}<0\);
the symmetric argument proves it halts exactly when \(x<0\).
Run both tests while reading the same cached prefix to
recognize \(x\ne0\).
On that promised domain exactly one kind of strict
certificate eventually appears, so output a bit for its
sign and pad the tail indefinitely. The finite decoding
of that bit as \(-1\) or \(1\) gives the usual sign value.
At zero neither certificate appears, as required.
\end{proof}
\end{annexitem}

\begin{annexitem}{annex:p04-rs-continuous-noncomputable}{thm:p04-rs-continuous-noncomputable}{Noncomputable constants despite continuity}
\begin{proof}
There are countably many finite programs, so at most
countably many productive computable names and hence
countably many \(\delta_{\mathrm C}\)-computable reals.
There are uncountably many reals: binary streams are
uncountable by diagonalization, and the map
\[
 b\longmapsto\sum_{n=0}^{\infty}2b(n)3^{-n-1}
\]
injects them into \(\mathbb R\).
Indeed, for two distinct streams whose first difference
is at \(k\), the differing term has magnitude
\(2\cdot3^{-k-1}\), while the maximum sum of all later
differences is \(3^{-k-1}\), strictly smaller.
Thus the two real values differ.
Choose a real \(a\) with no computable Cauchy name.

The map \(f(x)=a\) is continuous, since every open-set
preimage is either the whole domain or empty.
Choose any valid name \(r\) of \(a\).
The constant operator \(H(p)=r\) is continuous on all
of \(\mathcal B\), and realizes \(f\).
It is not claimed to have a finite program.
If a machine realizer of \(f\) existed, feed it the
computable constant zero-code name of zero.
The resulting productive output would be a computable
name of \(a\), a contradiction.
Hence this continuous and continuously realizable
total real map is not computable.
\end{proof}
\end{annexitem}
